\documentclass[11pt]{article}
\usepackage[T1]{fontenc}
\usepackage{lmodern,amsmath,amssymb,amsthm,mathtools}
\usepackage[margin=1in]{geometry}
\usepackage{microtype,tikz-cd}
\usepackage[colorlinks=true,linkcolor=blue,citecolor=blue,urlcolor=blue]{hyperref}
\hypersetup{pdftitle={A recursion category for Blum--Shub--Smale computation},pdfauthor={Florian Lengyel}}
\newtheorem{theorem}{Theorem}[section]
\newtheorem{proposition}[theorem]{Proposition}
\newtheorem{lemma}[theorem]{Lemma}
\newcommand{\R}{\mathbb R}
\newcommand{\N}{\omega}
\newcommand{\C}{\mathcal C_{\mathrm{BSS}}}
\newcommand{\D}{\mathcal D_{\mathrm{BSS}}}
\newcommand{\pfn}{\rightharpoonup}
\newcommand{\dom}{\operatorname{dom}}
\newcommand{\enc}{\operatorname{enc}}
\newcommand{\Split}{\operatorname{Split}_{\mathrm r}}
\title{A recursion category for Blum--Shub--Smale computation}
\author{Florian Lengyel\thanks{Email:
\href{mailto:flengyel@gradcenter.cuny.edu}{\texttt{flengyel@gradcenter.cuny.edu}}}}
\date{October 22, 2026}

\begin{document}
\maketitle
\begin{abstract}
We construct a dominical recursion category whose arrows are
Blum--Shub--Smale (BSS) computable partial functions over the real numbers.
The objects are positive finite powers of a space of tagged finite real
tuples. An explicit computable bijection identifies this space with its
square, and program specialization makes a universal evaluator a Turing
morphism. The category of BSS-computable partial functions between
finite-dimensional real spaces embeds fully faithfully in the restriction
idempotent completion.
\end{abstract}

The existence of a recursion-category presentation of BSS computation
was stated in \cite[\S3, p.~25]{Lengyel01}, with the construction omitted.
We give the construction over $\R$. The computation model and its
universal machines are due to Blum, Shub, and Smale
\cite{BSS89,BCSS98}; the categorical definitions are those of Di~Paola
and Heller \cite{DPH87} and Heller \cite{Heller90}.
Vigna \cite{Vigna96} gives another categorical characterization of
BSS computability over commutative ordered rings.
Our construction specifies an isotype and a Turing morphism and
identifies the fixed-dimensional computations in its completion.

\section{Finite data and a computable pairing}
\label{sec:data}

We use deterministic BSS machines over the ordered field $\R$, with
exact arithmetic, equality and order tests, and finite programs.
Each program may contain finitely many arbitrary real constants.
Running time is unrestricted. Inputs and outputs are finite tuples
whose lengths are recorded; machines have unbounded tape storage.
These are the finite-sequence conventions used for universal
computation in \cite[\S8]{BSS89} and \cite[Chapter~3]{BCSS98}.
Division by zero is treated as divergence. A \emph{partial function}
$f:A\pfn B$ is a function from a subset $\dom(f)\subseteq A$ to $B$;
the subset $\dom(f)$ is its \emph{domain of definition}.
The partial function is \emph{total} when $\dom(f)=A$.
It is BSS-computable when a machine halts exactly on this subset
and returns $f(a)$ on each input $a\in\dom(f)$.

Put $\N=\{0,1,\ldots\}$ and
$\R^{<\omega}=\coprod_{n\in\N}\R^n$, including the empty tuple.
For a tuple $a$, write $|a|$ for its length and $a^\frown b$ for
concatenation. Define
\[
 X=\N\times\R^{<\omega}.
\]
An element $(k,a)$, with $a=(a_0,\ldots,a_{n-1})$, is represented by
the finite record
\[
 \enc(k,a)=(n,k,a_0,\ldots,a_{n-1}).
\]
The two initial entries are natural numbers stored as real numbers.
An entry can be tested for membership in $\N$ by rejecting negative
values and repeatedly subtracting $1$ until it lies in $[0,1)$;
the original entry is a natural number exactly when the remainder is $0$.
For each fixed $r\geq1$, an element of $X^r$ is represented by the
concatenation of its $r$ records. The length fields determine their
boundaries. Computability of partial maps between these sets always
refers to these representations.

The natural-number tag will record how a tuple was split.
Define the Cantor pairing $c:\N^2\to\N$ by
$c(k,l)=(k+l)(k+l+1)/2+l$. Its inverse is computable by a finite
search through successive diagonals of $\N^2$. Such searches, and
integer quotient and remainder by a positive integer, are implementable
by BSS loops using addition, subtraction, and comparisons.

\begin{lemma}\label{lem:pair}
The map $b:X\times X\to X$ defined, for $|a|=n$ and $|a'|=m$, by
\[
 b\bigl((k,a),(l,a')\bigr)
   =\bigl((n+m+1)c(k,l)+n,\ a^\frown a'\bigr)
\]
is a bijection, and both $b$ and $b^{-1}$ are total BSS-computable.
\end{lemma}

\begin{proof}
Given $(q,z)\in X$, put $d=|z|$ and compute the unique quotient and
remainder $q=(d+1)t+n$, where $0\leq n\leq d$.
Recover $(k,l)=c^{-1}(t)$, split $z$ after its first $n$ entries,
and write the resulting tuples as $a,a'$. Then
$b^{-1}(q,z)=((k,a),(l,a'))$.
These operations are inverse to the displayed formula, including
when either tuple is empty. Every loop terminates on the indicated
natural-number data, and the remaining operations copy finite lists.
\end{proof}

Iterating $b$ gives chosen total computable bijections
$b_r:X^r\to X$ for all $r\geq1$, with $b_1=1_X$.
The tag is independent of the tuple length; this independence allows
every output tag to occur in Lemma~\ref{lem:pair}.

\section{The category of partial computations}
\label{sec:category}

Define $\C$ to have objects $X^r$, $r\geq1$, and all BSS-computable
partial functions $f:X^r\pfn X^s$ as arrows. Equality of arrows means
equality of domains and of values on those domains. Composition
$g\circ f$ applies $f$ first and is defined exactly where $f$ is
defined and $g$ is defined at its value.

A \emph{partial identity} on a subset $S\subseteq A$ is the partial
function $A\pfn A$ with domain $S$ and value $a$ at each $a\in S$.
The \emph{restriction} $\bar f:A\pfn A$ of $f:A\pfn B$ is the
partial identity on $\dom(f)$. Thus $f$ is total exactly when
$\bar f=1_A$, where $1_A$ is the identity
function on $A$. Write $0_{A,B}:A\pfn B$ for the nowhere-defined
partial function, whose domain is empty; subscripts are omitted
when the source and target are determined by context.
For $f:A\pfn B$ and $g:C\pfn D$, their \emph{near-product} is
the partial function
$f\boxtimes g:A\times C\pfn B\times D$ with value
$(f(a),g(c))$ and domain $\dom(f)\times\dom(g)$.
On the underlying sets, cartesian products are identified by
concatenating the tuples of $X$-coordinates; thus the near-product
of objects is $X^r\boxtimes X^s=X^{r+s}$.
The projections $(a,b)\mapsto a$ and $(a,b)\mapsto b$ and the
diagonal $a\mapsto(a,a)$ are total.

A \emph{$P$-category} is a category $\mathcal K$ with a functor
$\boxtimes:\mathcal K\times\mathcal K\to\mathcal K$, called the
near-product, a natural family $\Delta_A:A\to A\boxtimes A$,
and projections $\pi_0:A\boxtimes B\to A$ and
$\pi_1:A\boxtimes B\to B$, subject to the following axioms
\cite[\S5]{Heller90}. For maps $f:T\to A$ and $g:T\to B$, write
$\langle f,g\rangle=(f\boxtimes g)\circ\Delta_T$.
Object subscripts on projections are suppressed when their types
are determined by the formula.

First, $\pi_i\circ\Delta_A=1_A$ for $i=0,1$, and
$\langle\pi_0,\pi_1\rangle=1_{A\boxtimes B}$.
Second, let $\mathcal P$ be the smallest subcategory containing all
objects, identities, and projections and closed under $\boxtimes$.
For every $f:A\to B$ in $\mathcal K$ and $p:C\to D$ in
$\mathcal P$, the projections satisfy
\[
 \pi_0\circ(f\boxtimes p)=f\circ\pi_0,
 \qquad
 \pi_1\circ(p\boxtimes f)=f\circ\pi_1.
\]
Third, the following families are required to be natural:
\[
\begin{aligned}
 \alpha_{A,B,C}
  &=\langle\pi_0\circ\pi_0,
          \langle\pi_1\circ\pi_0,\pi_1\rangle\rangle
    :(A\boxtimes B)\boxtimes C\longrightarrow A\boxtimes(B\boxtimes C),\\
 \sigma_{A,B}
  &=\langle\pi_1,\pi_0\rangle
    :A\boxtimes B\longrightarrow B\boxtimes A.
\end{aligned}
\]
For an arrow $f:A\to B$ in a $P$-category, define
its restriction by $\bar f=\pi_0\circ\langle1_A,f\rangle$;
the arrow is total when $\bar f=1_A$.
In the category of sets and partial functions this agrees with the
partial-identity definition above. In this category,
$\alpha$ reassociates triples and $\sigma$ interchanges the two
coordinates. The stated axioms hold for the near-product, diagonal,
and projections defined above. In particular, every map in
$\mathcal P$ is total, which gives the two projection equations.
A subcategory closed under this structure inherits the axioms.

A family of \emph{zero arrows} consists of one arrow $0_{A,B}$ for
each pair of objects, absorbing composition and near-products:
every composite or near-product with a zero arrow is the zero arrow
of its type. The nowhere-defined partial functions have this property.
A $P$-category with such zero arrows is \emph{prodominical}.
A map $f$ is \emph{weakly total} when
$f\circ h=0$ implies $h=0$ for every composable $h$.
By \cite[Proposition~5.1, p.~1261]{Heller90}, a prodominical category
is \emph{dominical} if every weakly total map is total and
$f\boxtimes g=0$ implies $f=0$ or $g=0$.
An \emph{isotype} is a category in which all objects are isomorphic.

\begin{proposition}\label{prop:dominical}
The category $\C$ is a dominical $P$-category and an isotype.
\end{proposition}

\begin{proof}
Identities and composition are computable: for a composite, run the
first machine and then the second. To compute $\bar f$, save the
input, run a machine for $f$, and return the saved input if it halts.
To compute $f\boxtimes g$, run the two machines successively on
their respective inputs and return both outputs. This halts exactly
on $\dom(f)\times\dom(g)$. Projections, diagonals, and the finite
coordinate rearrangements are computable by copying records.
Thus the $P$-category structure is inherited from partial functions.

The nowhere-defined map $0_{A,B}:A\pfn B$ is computed by a
nonterminating loop. It absorbs composition and near-products,
so $\C$ is prodominical.

Suppose that $f:A\pfn B$ is not total, and choose
$a\in A\setminus\dom(f)$. The constant map $h:X\to A$ with value
$a$ is total and BSS-computable: its program uses the finitely many
real entries of $a$ as constants. Since $X$ is nonempty,
$h\ne0$, while $f\circ h=0$. Hence $f$ is not weakly total.
Also, if $f\ne0$ and $g\ne0$, choose one point in each domain;
their pair belongs to $\dom(f\boxtimes g)$. Heller's two
conditions for dominicality follow.

Finally, the maps $b_r$ and their inverses are arrows of $\C$.
They identify every object with $X$.
\end{proof}

\section{Universal evaluation and specialization}
\label{sec:universal}

In a prodominical isotype, a \emph{Turing morphism} is an arrow
$\tau:W\boxtimes A\pfn B$ such that for every object $V$ and
every arrow $\varphi:V\boxtimes A\pfn B$ there is a total arrow
$s:V\to W$ satisfying
$\varphi=\tau\circ(s\boxtimes1_A)$.
A \emph{recursion category} is a prodominical isotype with a Turing
morphism \cite[\S9, pp.~1264--1265]{Heller90}.
The equality in this definition is equality of partial maps.

We describe the program representation used for the evaluator.
A BSS program can be written as a finite instruction list over a fixed
set of instruction types: arithmetic, copying, tape movement,
conditional branching, and halting. Polynomial and rational
instructions are expanded into finite arithmetic instruction lists.
The instruction types, register addresses, branch destinations, and
list lengths are natural numbers. The real constants are stored
separately in a finite table and referred to by their indices.
For a finite list $(u_0,\ldots,u_{L-1})$ of natural numbers, put
$h_L=0$, define $h_i=c(u_i,h_{i+1})$ recursively backwards for $0\leq i<L$,
and define its code to be $c(L,h_0)$. To decode, recover $L$ and
$h_0$, unpair $L$ times, and require the final remainder to be $0$.
Apply this encoding to the integer instruction data. Together
with the real constant table, it gives a program code in $X$.

An interpreter decodes the instruction list, maintains the simulated
tape and instruction pointer, and executes one instruction at a time.
Integer addresses are handled by finite scans, and a constant-table
lookup copies the indicated real entry. Thus one interpreter performs
the indicated computation without placing the simulated constants in
its own finite program. Invalid instructions or indices cause it to
loop. An output is accepted only if it is a valid record for an
element of $X$. This gives a BSS-computable partial evaluator
$U:X\boxtimes X\pfn X$. It is the usual universal-machine
construction of \cite[\S8, pp.~34--36]{BSS89}, with instruction
data and real constants stored in separate fields.

\begin{lemma}\label{lem:specialization}
For every $r\geq1$ and every BSS-computable partial function
$\varphi:X^r\times X\pfn X$, there is a total BSS-computable
$s_\varphi:X^r\to X$ such that, for every $v\in X^r$ and $x\in X$,
$U(s_\varphi(v),x)$ and $\varphi(v,x)$ have the same definedness
and, when defined, the same value.
\end{lemma}

\begin{proof}
Fix a program $P$ computing $\varphi$ on the concatenated input
records. Given $v\in X^r$, construct a program whose input is $x$.
Its initialization instructions save $x$, write the records of $v$
into the parameter part of the simulated input, and append the
saved record of $x$. It then runs $P$ on that input.

To produce this program, copy the finitely many real entries of $v$
to its constant table, together with the constants of $P$.
Generate one finite initialization prefix of the required length,
then append the fixed instruction list for $P$, adjusting addresses
and branch destinations. All loops in this construction are bounded
by the lengths of the input records or of the generated instruction
list. The construction never runs $P$. It therefore defines a total
BSS-computable code map $s_\varphi$.

The generated program on input $x$ runs $P$ on $(v,x)$ after a
finite initialization. Consequently it halts exactly when
$\varphi(v,x)$ is defined and returns the same value.
\end{proof}

\begin{samepage}
The following diagram expresses the specialization equation.
Its top arrow is total, and the two paths have the same restriction
domain by Lemma~\ref{lem:specialization}.
\[
\begin{tikzcd}[column sep=large,row sep=large]
 X^r\boxtimes X
   \arrow[r,"s_\varphi\boxtimes1_X"]
   \arrow[dr,"\varphi"']
   &X\boxtimes X\arrow[d,"U"]\\
   &X .
\end{tikzcd}
\]
\end{samepage}
\begin{theorem}\label{thm:main}
Let $X=\omega\times\coprod_{n\geq0}\R^n$, with the record
representations of Section~\ref{sec:data}. The category with objects
$X^r$, $r\geq1$, and all unrestricted-time BSS-computable partial
functions between them, allowing finitely many arbitrary real
constants in each program, is a dominical recursion category.
\end{theorem}

\begin{proof}
Proposition~\ref{prop:dominical} makes it a dominical isotype.
Every object $V$ is $X^r$ for some $r\geq1$, so
Lemma~\ref{lem:specialization} says exactly that
$U:X\boxtimes X\pfn X$ is a Turing morphism.
\end{proof}

\section{Recovery of fixed-dimensional computations}
\label{sec:typed}

Let $\D$ be the category with objects $\R^n$, $n\geq0$, and all
BSS-computable partial functions between them. A \emph{restriction
idempotent} is an arrow $e:A\pfn A$ satisfying $e=\bar e$;
for partial functions, these are precisely the partial identities.
For a category of partial maps $\mathcal C$ closed under restrictions,
its \emph{restriction idempotent completion} $\Split(\mathcal C)$
has objects $(A,e)$, where $e:A\pfn A$ is a restriction idempotent
belonging to $\mathcal C$.
An arrow $h:(A,e)\pfn(B,d)$ is an arrow of $\mathcal C$ satisfying
$h=d\circ h\circ e$. Its composition is inherited, and the identity
of $(A,e)$ is $e$. This is the splitting construction for restriction
idempotents discussed in \cite[\S2.1, p.~187]{CockettHofstra08}.
Restrictions are inherited from $\mathcal C$. If $\mathcal C$ has
the near-product of partial functions, define
$(A,e)\boxtimes(B,d)=(A\boxtimes B,e\boxtimes d)$ and use
the inherited near-product on arrows.

For $n\geq0$, let $S_n=\{(0,a)\in X:|a|=n\}$, and let
$e_n:X\pfn X$ be the partial identity on $S_n$.
The tag and length tests decide membership in $S_n$, so $e_n$
is BSS-computable. Define the total encoder
$i_n:\R^n\to X$ by $i_n(a)=(0,a)$ and the partial decoder
$d_n:X\pfn\R^n$ by $d_n(0,a)=a$ on $S_n$.
They are BSS-computable and satisfy
$d_n\circ i_n=1_{\R^n}$ and $i_n\circ d_n=e_n$.
Their endpoint types here refer to sets and partial functions.

\begin{proposition}\label{prop:embedding}
There is a fully faithful functor
$J:\D\to\Split(\C)$ given by
\[
 J(\R^n)=(X,e_n),\qquad
 J(f)=i_m\circ f\circ d_n
 \quad(f:\R^n\pfn\R^m).
\]
\end{proposition}

\begin{proof}
The map $J(f)$ is BSS-computable, has domain contained in $S_n$,
and takes its values in $S_m$. Thus
$J(f)=e_m\circ J(f)\circ e_n$.
Moreover, $J(1_{\R^n})=e_n$, which is the identity of $(X,e_n)$.
For composable $f,g$, the identity $d_m\circ i_m=1_{\R^m}$
gives $J(g)\circ J(f)=J(g\circ f)$, including equality of domains.

Conversely, an arrow $h:(X,e_n)\pfn(X,e_m)$ decodes to the
BSS-computable partial function $d_m\circ h\circ i_n$.
Encoding and decoding are inverse on these hom-sets. In particular,
$i_m\circ(d_m\circ h\circ i_n)\circ d_n
=e_m\circ h\circ e_n=h$, and the other composite is the identity
by $d_j\circ i_j=1_{\R^j}$.
\end{proof}

This identifies $\D$ with the full subcategory on the objects
$(X,e_n)$. Moreover,
$\overline{J(f)}=J(\bar f)$, since both sides are the partial
identity on $i_n(\dom(f))$.
The representation also respects the near-products up to
computable isomorphism: on the supported domains, concatenate the
real tuples to obtain $S_n\times S_m\cong S_{n+m}$, and split after
the first $n$ entries for the inverse. Thus the fixed-dimensional
BSS computations are recovered with their partial-map structure.

\begin{thebibliography}{9}
\small
\setlength{\itemsep}{3pt}

\bibitem{BSS89}
Lenore Blum, Mike Shub, and Steve Smale,
On a theory of computation and complexity over the real numbers:
$NP$-completeness, recursive functions and universal machines,
\emph{Bulletin of the American Mathematical Society (N.S.)}
\textbf{21} (1989), no.~1, 1--46.
\href{https://doi.org/10.1090/S0273-0979-1989-15750-9}{doi:10.1090/S0273-0979-1989-15750-9}.

\bibitem{BCSS98}
Lenore Blum, Felipe Cucker, Michael Shub, and Steve Smale,
\emph{Complexity and Real Computation},
Springer, New York, 1998.
\href{https://doi.org/10.1007/978-1-4612-0701-6}{doi:10.1007/978-1-4612-0701-6}.

\bibitem{CockettHofstra08}
J.~Robin~B. Cockett and Pieter~J.~W. Hofstra,
Introduction to Turing categories,
\emph{Annals of Pure and Applied Logic}
\textbf{156} (2008), nos.~2--3, 183--209.
\href{https://doi.org/10.1016/j.apal.2008.04.005}{doi:10.1016/j.apal.2008.04.005}.

\bibitem{DPH87}
Robert~A. Di~Paola and Alex Heller,
Dominical categories: recursion theory without elements,
\emph{Journal of Symbolic Logic} \textbf{52} (1987), no.~3, 594--635.
\href{https://doi.org/10.2307/2274352}{doi:10.2307/2274352}.

\bibitem{Heller90}
Alex Heller,
An existence theorem for recursion categories,
\emph{Journal of Symbolic Logic} \textbf{55} (1990), no.~3, 1252--1268.
\href{https://doi.org/10.2307/2274486}{doi:10.2307/2274486}.

\bibitem{Lengyel01}
Florian Lengyel,
\emph{Recursion categories of coalgebras},
preprint, 2001.
\href{https://arxiv.org/abs/math/0105256}{arXiv:math/0105256}.

\bibitem{Vigna96}
Sebastiano Vigna,
On the relations between distributive computability and the BSS model,
\emph{Theoretical Computer Science} \textbf{162} (1996), no.~1, 5--21.
\href{https://doi.org/10.1016/0304-3975(95)00247-2}{doi:10.1016/0304-3975(95)00247-2}.

\end{thebibliography}
\end{document}
